
\newcolumntype{L}[1]{>{\raggedright\arraybackslash}p{#1}}
\newcolumntype{C}[1]{>{\centering\arraybackslash}p{#1}}

\begin{tcolorbox}[colback=primaryColor,colframe=primaryColor,arc=3mm,halign=center,valign=center,height=2.5cm]
 \color{white}\LARGE\bfseries INSPECTION CARD \\[0.2cm]
 \color{white!80}\large UX Design Project: Help People Help People --- TPER S.p.A. \\[0.1cm]
 \color{white!60}\normalsize Phase D --- Cognitive Walkthrough
\end{tcolorbox}

\vspace{0.3cm}

% =====================================================================
% INSPECTION DATA (aligned to template)
% =====================================================================
\section{Inspection Data}

\begin{tcolorbox}[colback=backgroundColor,colframe=borderColor,arc=2mm]
\small

\noindent\textbf{Name of version:} TPER Redesign --- v1.0 (Phase C -- Design Proposal)

\vspace{0.3cm}
\noindent\textbf{Inspection context data:}
\begin{itemize}
 \item \textbf{Date:} June 15--16, 2026
 \item \textbf{Team members:} 2 evaluators (design team members)
 \item \textbf{Roles:} Evaluator 1 (narrator, guides the simulation), Evaluator 2 (listener, records critical issues)
 \item \textbf{Mode:} Individual evaluation on static prototype (HTML/CSS), followed by comparison and consensus
\end{itemize}

\vspace{0.3cm}
\noindent\textbf{Inspection method:}
\begin{itemize}
 \item[$\bullet$] \textbf{Cognitive Walkthrough} (selected)
 \item[$\circ$] Formal Action Analysis
 \item[$\circ$] Informal Action Analysis
 \item[$\circ$] Heuristic Analysis
\end{itemize}

\vspace{0.3cm}
\noindent\textbf{Output of inspection:} The 5 analyzed tasks are documented below. For each step, the 4 CW questions (Goal, Visibility, Mapping, Feedback) are evaluated. A qualitative justification is provided for each critical issue detected.

\end{tcolorbox}

\vspace{0.3cm}

% =====================================================================
% INTRODUCTION TO COGNITIVE WALKTHROUGH
% =====================================================================
\section{Introduction to the Cognitive Walkthrough}

\begin{tcolorbox}[colback=backgroundColor,colframe=borderColor,arc=2mm]
\small
The Cognitive Walkthrough is an inspection method that evaluates learnability through simulated task execution. For each step, 4 questions are answered (Goal, Visibility, Mapping, Feedback). Evaluators simulate the profiles of the personas defined in Phase B, assuming minimal digital literacy and a domestic context without assistance.

\vspace{0.2cm}
\noindent\textbf{Selected tasks:}
\begin{enumerate}
 \item Subscription Renewal (Guest Flow)
 \item Subscription Purchase (Login Flow)
 \item Discovering Concessions
 \item Purchase Guide (Where to Buy)
 \item Card Verification and Top-up
\end{enumerate}
\end{tcolorbox}

\vspace{0.5cm}

% =====================================================================
% TASK 1: SUBSCRIPTION RENEWAL (GUEST)
% =====================================================================
\section{Task 1: Subscription Renewal --- Guest Flow}

\begin{tcolorbox}[colback=primaryColor!10,colframe=primaryColor,arc=2mm]
\small
\textbf{User profile:} Liliana (70, low digital literacy) needs to renew her monthly subscription that expires tomorrow. She tries to do it on her own, without assistance.

\textbf{Starting point:} Homepage \texttt{tper.it}. \textbf{Goal:} Complete the renewal without registration.
\end{tcolorbox}

\vspace{0.3cm}

\subsection{Step 1: Finding the ``Where to Buy'' section}

\begin{tcolorbox}[colback=backgroundColor,colframe=borderColor,arc=2mm]
\small
\begin{tabular}{|L{3cm}|L{11cm}|}
\hline
\textbf{User action} & \multicolumn{1}{L{11cm}|}{From the homepage, clicks ``Where to Buy'' in the main navigation or on the ``Where to Buy'' card in the ``What do you want to do?'' grid} \\
\hline
\textbf{Q1 -- Goal} & The user wants to renew a subscription. The term ``Where to Buy'' may not be immediately associated with renewal. \\
\hline
\textbf{Q2 -- Visibility} & The ``Where to Buy'' item is clearly visible in the main menu (position \#4) and as a card on the homepage (green card with cart icon). \\
\hline
\textbf{Q3 -- Mapping} & ``Where to Buy'' suggests purchase, not renewal. The user may hesitate. However, the ``Subscriptions'' card (purple card) is closer to the concept of ``subscription renewal''. \\
\hline
\textbf{Q4 -- Feedback} & Navigation is a simple page transition, no specific feedback. \\
\hline
\end{tabular}

\vspace{0.3cm}
\noindent\textbf{Critical issue:} Q3 is critical: the term ``Where to Buy'' is ambiguous for a renewal operation. The user might click ``Subscriptions'' instead of ``Where to Buy'', moving away from the guided channel selection flow. Although ``Subscriptions'' provides fare information, it does not lead directly to the renewal channel, requiring an additional navigation step.
\end{tcolorbox}

\vspace{0.3cm}

\subsection{Steps 2--3: Using the interactive guide}

\begin{tcolorbox}[colback=backgroundColor,colframe=borderColor,arc=2mm]
\small
\noindent The user clicks ``We'll help you choose'' and answers the 3 questions (service, channel, payment). No critical issues: the guide is clear, options are large and clickable, the ``Question X of 3'' progress bar and the ``Back'' button provide adequate feedback. The wizard takes the user to the result ``solweb.tper.it'' with a direct link.
\end{tcolorbox}

\vspace{0.3cm}

\subsection{Step 4: Reaching solweb and choosing ``Impersonal Renewal''}

\begin{tcolorbox}[colback=backgroundColor,colframe=borderColor,arc=2mm]
\small
\begin{tabular}{|L{3cm}|L{11cm}|}
\hline
\textbf{User action} & \multicolumn{1}{L{11cm}|}{From the guide, gets ``solweb.tper.it'' as the result. Clicks ``Go to solweb'' which leads to \texttt{login.html}. Here, clicks ``Impersonal Renewal'' among the 4 Quick Actions, arriving at \texttt{rinnova.html}. Chooses ``Impersonal Subscription'' and then ``Renew without Login''.} \\
\hline
\textbf{Q1 -- Goal} & The user wants to renew without creating an account. ``Impersonal Renewal'' and ``Renew without Login'' match this need. \\
\hline
\textbf{Q2 -- Visibility} & The ``Quick Actions'' are in a clearly visible section under the login form, with large buttons. The decision screen on \texttt{rinnova.html} shows two side-by-side cards with clear icons and descriptions. \\
\hline
\textbf{Q3 -- Mapping} & ``Impersonal Subscription'' $\rightarrow$ ``Non-nominative. You just need your card number'' $\rightarrow$ ``Renew without Login''. The language is simple and descriptive. \\
\hline
\textbf{Q4 -- Feedback} & The progress bar shows ``Step 1 of 2''. The breadcrumb shows ``Login / Impersonal Renewal''. Good orientation. \\
\hline
\end{tabular}

\vspace{0.3cm}
\noindent The guest flow is clear, options are well explained. The user has a direct path for renewal without needing to create an account.

\vspace{0.2cm}
\noindent\textbf{Critical issue:} \textbf{CW-08 (Medium, Q2).} On \texttt{login.html}, the Quick Actions are at the bottom of the page, below the login form that visually dominates the screen. For Liliana, the login form is the focal point: she risks not noticing the guest options, attempting to log in without credentials, and abandoning the task.
\end{tcolorbox}

\vspace{0.3cm}

\subsection{Step 5: Filling in the renewal form}

\begin{tcolorbox}[colback=backgroundColor,colframe=borderColor,arc=2mm]
\small
\begin{tabular}{|L{3cm}|L{11cm}|}
\hline
\textbf{User action} & \multicolumn{1}{L{11cm}|}{Enters the card number (format TPER-XXXXXX) and clicks ``Confirm and Proceed''} \\
\hline
\textbf{Q1 -- Goal} & The user understands they must enter the card number to identify the subscription to be renewed. \\
\hline
\textbf{Q2 -- Visibility} & The ``Card Number / Subscription Code'' field is the only field on the form. Clearly visible with label and placeholder. \\
\hline
\textbf{Q3 -- Mapping} & ``Card number'' $\rightarrow$ ``identify my subscription'' is a clear mapping. Many users are familiar with this type of form (e.g., phone top-ups). \\
\hline
\textbf{Q4 -- Feedback} & Validation is client-side: if the field is empty, it shows an inline error (\texttt{showFieldError}). On submit, redirect to \texttt{pagamento.html}. Quick feedback. \\
\hline
\end{tabular}

\vspace{0.3cm}
\noindent Minimal and clear form.

\vspace{0.2cm}
\noindent\textbf{Issue identified:} The format ``TPER-XXXXXX'' has no formal client-side validation. A user who enters an incorrect format (e.g., numbers only) will still be redirected to payment, only to discover the error later. It is recommended to add a validation regex with immediate feedback.
\end{tcolorbox}

\vspace{0.3cm}


\vspace{0.3cm}

% =====================================================================
% TASK 2: SUBSCRIPTION PURCHASE (LOGIN FLOW)
% =====================================================================
\section{Task 2: Subscription Purchase --- Login Flow}

\begin{tcolorbox}[colback=primaryColor!10,colframe=primaryColor,arc=2mm]
\small
\textbf{User profile:} David (34, English speaker, high digital literacy) wants to buy an annual subscription. He already has a solweb account. He tries to do everything on his own.

\textbf{Starting point:} ``Where to Buy'' page on \texttt{tper.it}. \textbf{Goal:} Complete the purchase through to confirmation.
\end{tcolorbox}

\vspace{0.3cm}

\subsection{Step 1: Reaching login and signing in}

\begin{tcolorbox}[colback=backgroundColor,colframe=borderColor,arc=2mm]
\small
\begin{tabular}{|L{3cm}|L{11cm}|}
\hline
\textbf{User action} & \multicolumn{1}{L{11cm}|}{From ``Where to Buy'', clicks ``Go to solweb.tper.it'' $\rightarrow$ \texttt{login.html}. Enters email/CF and password, then clicks ``Sign in''.} \\
\hline
\textbf{Q1 -- Goal} & The user knows they must sign in to buy a subscription. The informational banner ``To purchase a subscription, you must sign in'' confirms this. \\
\hline
\textbf{Q2 -- Visibility} & The login form is the main content of the page. The ``Fiscal Code / Email'' and ``Password'' fields are clearly visible. \\
\hline
\textbf{Q3 -- Mapping} & ``Enter your credentials'' $\rightarrow$ ``sign in to your account'' is the universal login mapping. The user recognizes it. However, for an English-speaking user, ``Fiscal Code'' is an unfamiliar Italian concept. \\
\hline
\textbf{Q4 -- Feedback} & On submit, immediate redirect to \texttt{dashboard.html?login=ok}. Clear feedback of successful login. In case of error, an error message with explanation and alternatives. \\
\hline
\end{tabular}

\vspace{0.3cm}
\noindent\textbf{Critical issue:} Q3 is critical: ``Fiscal Code'' is not an international concept. The English-speaking user may not understand what to enter. The placeholder ``or Email'' only partially mitigates the problem: the user must still pause, read, and interpret before finding the alternative.
\end{tcolorbox}

\vspace{0.3cm}

\subsection{Step 2: Dashboard and plan selection}

\begin{tcolorbox}[colback=backgroundColor,colframe=borderColor,arc=2mm]
\small
\begin{tabular}{|L{3cm}|L{11cm}|}
\hline
\textbf{User action} & \multicolumn{1}{L{11cm}|}{On the dashboard (2 cards: ``Purchase'', ``Personal Renewal''), clicks ``Purchase'' $\rightarrow$ \texttt{piani.html} with 3 plans: Monthly (51\euro), Annual (390\euro), Suburban (from 35\euro/month). Chooses Annual and clicks ``Buy Online''.} \\
\hline
\textbf{Q1 -- Goal} & The user wants to buy. ``Purchase'' on the dashboard is explicit. The plans show clear prices and descriptions. \\
\hline
\textbf{Q2 -- Visibility} & The plan cards are large, with highlighted prices and ``Select'' buttons. The progress bar shows ``Step 2 of 4''. \\
\hline
\textbf{Q3 -- Mapping} & ``Select'' on the annual plan $\rightarrow$ ``I want this subscription'' is direct and unambiguous. \\
\hline
\textbf{Q4 -- Feedback} & The progress bar updates, the breadcrumb shows the path. On click, redirect to \texttt{documenti.html}. \\
\hline
\end{tabular}

\vspace{0.3cm}
\noindent Linear flow, visible prices, clear selection.
\end{tcolorbox}

\vspace{0.3cm}

\subsection{Step 3: Document upload}

\begin{tcolorbox}[colback=backgroundColor,colframe=borderColor,arc=2mm]
\small
\begin{tabular}{|L{3cm}|L{11cm}|}
\hline
\textbf{User action} & \multicolumn{1}{L{11cm}|}{On \texttt{documenti.html}, sees the required documents checklist (ID, fiscal code, ID photo). Confirms they have them and clicks ``Proceed to Payment''.} \\
\hline
\textbf{Q1 -- Goal} & The user understands that some documents are needed for the annual subscription. The checklist shows them. \\
\hline
\textbf{Q2 -- Visibility} & The checklist is well structured with icons, descriptions, and notes. \\
\hline
\textbf{Q3 -- Mapping} & ``Do I have these documents?'' $\rightarrow$ ``Confirm and proceed''. The mapping is clear: the user verifies they have what is needed. \\
\hline
\textbf{Q4 -- Feedback} & Progress bar ``Step 3 of 4''. On submit, redirect to \texttt{pagamento.html}. Feedback OK. \\
\hline
\end{tabular}

\vspace{0.3cm}
\noindent The preventive checklist is a good solution to the ``arriving at the counter without documents'' problem.
\end{tcolorbox}

\vspace{0.3cm}

\subsection{Step 4: Payment and confirmation}

\begin{tcolorbox}[colback=backgroundColor,colframe=borderColor,arc=2mm]
\small
\begin{tabular}{|L{3cm}|L{11cm}|}
\hline
\textbf{User action} & \multicolumn{1}{L{11cm}|}{On \texttt{pagamento.html}, sees the order summary (annual plan 390\euro), enters payment details (card, expiry, CVV), clicks ``Pay 390\euro'' $\rightarrow$ redirect to \texttt{conferma.html} with summary and PDF receipt.} \\
\hline
\textbf{Q1 -- Goal} & The user wants to pay to complete the purchase. The summary and total are clearly visible. \\
\hline
\textbf{Q2 -- Visibility} & Complete payment form (cardholder, card number, expiry, CVV). ``Pay'' button with clearly visible amount. \\
\hline
\textbf{Q3 -- Mapping} & ``Fill in your card details'' $\rightarrow$ ``pay the amount shown'' is the e-commerce standard. Universal mapping. \\
\hline
\textbf{Q4 -- Feedback} & After payment, redirect to \texttt{conferma.html} with order summary, confirmation message, and link to download the PDF receipt. Excellent feedback. \\
\hline
\end{tabular}

\vspace{0.3cm}
\noindent Standard and reassuring payment flow.

\vspace{0.2cm}
\noindent\textbf{Issue identified:} \textbf{CW-03 (Medium).} The ``And now? What to do after purchase'' section is present at the bottom of the confirmation page, but it is not very visible: located below the order summary and the receipt download button, it requires the user to scroll below the fold. A user with low digital literacy may not notice the post-purchase instructions and board without a properly usable ticket.
\end{tcolorbox}

\vspace{0.3cm}


\vspace{0.3cm}

% =====================================================================
% TASK 3: DISCOVERING CONCESSIONS
% =====================================================================
\section{Task 3: Discovering Concessions}

\begin{tcolorbox}[colback=primaryColor!10,colframe=primaryColor,arc=2mm]
\small
\textbf{User profile:} Roberto (45, low motivation, low concentration) has heard he might be eligible for a discount on his subscription.

\textbf{Starting point:} ``Subscriptions'' page on \texttt{tper.it}. \textbf{Goal:} Find out if he is eligible for concessions and which documents are needed.
\end{tcolorbox}

\vspace{0.3cm}

\subsection{Step 1: Noticing the concessions banner}

\begin{tcolorbox}[colback=backgroundColor,colframe=borderColor,arc=2mm]
\small
\begin{tabular}{|L{3cm}|L{11cm}|}
\hline
\textbf{User action} & \multicolumn{1}{L{11cm}|}{On the ``Subscriptions'' page, sees a green banner with a trophy icon: ``Not sure if you're eligible for a discount? Discover in 30 seconds which concessions you can access'' with a ``Discover Concessions'' button.} \\
\hline
\textbf{Q1 -- Goal} & The user wants to know if they are eligible for discounts. The banner speaks exactly about this. \\
\hline
\textbf{Q2 -- Visibility} & The banner is at the top of the page, colored (green), with an appealing icon and short text. Very visible. \\
\hline
\textbf{Q3 -- Mapping} & ``Discover Concessions'' $\rightarrow$ ``find out if I am eligible for discounts'' is a direct and clear mapping. \\
\hline
\textbf{Q4 -- Feedback} & Click $\rightarrow$ navigation to \texttt{agevolazioni.html}. Standard page transition. \\
\hline
\end{tabular}

\vspace{0.3cm}
\noindent Well-designed banner: visible, targeted, with a clear CTA.
\end{tcolorbox}

\vspace{0.3cm}

\subsection{Step 2: Launching the wizard and answering questions}

\begin{tcolorbox}[colback=backgroundColor,colframe=borderColor,arc=2mm]
\small
\begin{tabular}{|L{3cm}|L{11cm}|}
\hline
\textbf{User action} & \multicolumn{1}{L{11cm}|}{On \texttt{agevolazioni.html}, clicks ``Discover My Concessions'' $\rightarrow$ answers 4 questions (age, ISEE, residence, travel frequency) by clicking on the options.} \\
\hline
\textbf{Q1 -- Goal} & The user understands that answering simple questions will lead to a personalized result. \\
\hline
\textbf{Q2 -- Visibility} & The questions are large, one at a time, with button options. The ``Question X of 4'' progress bar is visible. \\
\hline
\textbf{Q3 -- Mapping} & ``How old are you?'' $\rightarrow$ select the age range. ``Are you a resident of Bologna?'' $\rightarrow$ simple options. Questions in natural language, not bureaucratic. \\
\hline
\textbf{Q4 -- Feedback} & On clicking each answer, the next question appears immediately. ``Back'' button available. Progress bar updated. Excellent feedback. \\
\hline
\end{tabular}

\vspace{0.3cm}
\noindent Simple wizard, clear questions, non-technical language. The user does not need to know bureaucratic terms to answer.
\end{tcolorbox}

\vspace{0.3cm}

\subsection{Step 3: Reading the results}

\begin{tcolorbox}[colback=backgroundColor,colframe=borderColor,arc=2mm]
\small
\begin{tabular}{|L{3cm}|L{11cm}|}
\hline
\textbf{User action} & \multicolumn{1}{L{11cm}|}{Sees the results screen: ``We found X concessions for you!'' with a card for each concession, description, required documents, and a ``View Subscriptions'' link.} \\
\hline
\textbf{Q1 -- Goal} & The user wants to know if they are eligible for discounts. The results tell them clearly. \\
\hline
\textbf{Q2 -- Visibility} & The results are in colored cards with title, description, and icon. The required documents are listed. \\
\hline
\textbf{Q3 -- Mapping} & ``You are eligible for...'' $\rightarrow$ ``Here is what you need''. The connection is direct: the user understands what to do next. \\
\hline
\textbf{Q4 -- Feedback} & Success message with the number of concessions found. ``Restart'' button to redo the test. ``View Subscriptions'' link for further details. \\
\hline
\end{tabular}

\vspace{0.3cm}
\noindent Personalized, clear result, with an explicit call-to-action.

\vspace{0.2cm}
\noindent\textbf{Issue identified:} \textbf{CW-05 (Minor).} The result does not include a ``Print'' button to take the document list to the counter. For elderly or low-digital-literacy users, paper support is important.

\vspace{0.2cm}
\noindent\textbf{Issue:} \textbf{CW-10 (Medium, Q4).} If the wizard finds no concessions, the message ``No automatic concessions found'' may be perceived as failure. For Roberto (low motivation), this confirms the pointlessness of the attempt. A digital alternative that keeps the user in the flow is missing.
\end{tcolorbox}

\vspace{0.3cm}


\vspace{0.3cm}

% =====================================================================
% TASK 4: PURCHASE GUIDE (WHERE TO BUY)
% =====================================================================
\section{Task 4: Purchase Guide --- Where to Buy}

\begin{tcolorbox}[colback=primaryColor!10,colframe=primaryColor,arc=2mm]
\small
\textbf{User profile:} Fatima (68, Arabic speaker, illiterate in written Italian, minimal digital literacy) visits \texttt{tper.it} for the first time. She does not know the site is informational only. She is looking for a way to ``buy a bus ticket''.

\textbf{Starting point:} Homepage \texttt{tper.it}. \textbf{Goal:} Discover how and where to buy a ticket.
\end{tcolorbox}

\vspace{0.3cm}

\subsection{Step 1: Finding the ``Where to Buy'' section}

\begin{tcolorbox}[colback=backgroundColor,colframe=borderColor,arc=2mm]
\small
\begin{tabular}{|L{3cm}|L{11cm}|}
\hline
\textbf{User action} & \multicolumn{1}{L{11cm}|}{From the homepage, clicks ``Where to Buy'' (green card or navigation).} \\
\hline
\textbf{Q1 -- Goal} & Fatima wants to ``buy''. ``Where to Buy'' exactly matches her need. For a user with 1/5 literacy, the ``Where to Buy'' language is perfect. \\
\hline
\textbf{Q2 -- Visibility} & The ``Where to Buy'' card is the second in the ``What do you want to do?'' grid, clearly visible right after the hero section. Cart icon + clear title. \\
\hline
\textbf{Q3 -- Mapping} & ``Where to Buy'' $\rightarrow$ ``find out where to buy'' is literal and unambiguous. \\
\hline
\textbf{Q4 -- Feedback} & Navigation to \texttt{dove-comprare.html}. The page opens with the explanation that \texttt{tper.it} does not sell directly. \\
\hline
\end{tabular}

\vspace{0.3cm}
\noindent``Where to Buy'' is the best labeling choice for users with low literacy.
\end{tcolorbox}

\vspace{0.3cm}

\subsection{Step 2: Reading the informational banner}

\begin{tcolorbox}[colback=backgroundColor,colframe=borderColor,arc=2mm]
\small
\begin{tabular}{|L{3cm}|L{11cm}|}
\hline
\textbf{User action} & \multicolumn{1}{L{11cm}|}{Arriving on the page, reads the text: ``On tper.it it is not possible to purchase travel passes. This page guides you to the most suitable channel.''} \\
\hline
\textbf{Q1 -- Goal} & Fatima's goal (to buy) is not achievable here. But she must understand that the page will guide her elsewhere. \\
\hline
\textbf{Q2 -- Visibility} & The text is immediately below the h1, in a prominent position. \\
\hline
\textbf{Q3 -- Mapping} & ``It is not possible to purchase... but it guides you to the most suitable channel'' $\rightarrow$ ``I need to go somewhere else, but here I find out where''. The message is clear but may disappoint the user. \\
\hline
\textbf{Q4 -- Feedback} & The page immediately shows the ``We'll help you choose'' button and the channel cards. Fatima has visible options. \\
\hline
\end{tabular}

\vspace{0.3cm}
\noindent\textbf{Critical issue:} A user with 1/5 literacy might feel ``rejected'' by the ``you cannot buy here'' message. However, the interactive guide and cards offer an immediate alternative.

\vspace{0.2cm}
\noindent\textbf{Issue:} \textbf{CW-06 (Medium).} For a very elderly user, the concept of ``informational website'' is abstract. The banner explains that you cannot buy, but does not offer a one-click path for those who do not want to use the guide. It is recommended to add a quick path ``I want to buy now'' that leads directly to an operational channel.

\vspace{0.2cm}
\noindent\textbf{Issue:} \textbf{CW-09 (Medium, Q2).} The language selector uses the abbreviations IT, EN, AR in Latin script. For Fatima, who only reads Arabic, ``AR'' is not a recognizable signal. She may never discover that the site is available in Arabic.
\end{tcolorbox}

\subsection{Step 3: Using the interactive guide}

\begin{tcolorbox}[colback=backgroundColor,colframe=borderColor,arc=2mm]
\small
\begin{tabular}{|L{3cm}|L{11cm}|}
\hline
\textbf{User action} & \multicolumn{1}{L{11cm}|}{Clicks ``We'll help you choose'' $\rightarrow$ answers ``A ticket'' $\rightarrow$ ``No preference'' $\rightarrow$ ``Both / I don't know'' $\rightarrow$ Gets ``Roger App'' as the recommended channel.} \\
\hline
\textbf{Q1 -- Goal} & Fatima understands that the guide will take her to a channel where she can buy. \\
\hline
\textbf{Q2 -- Visibility} & The primary ``We'll help you choose'' button is clearly visible. The options are large. \\
\hline
\textbf{Q3 -- Mapping} & ``A ticket (single ride or carnet)'' $\rightarrow$ Fatima wants ``a bus ticket''. Direct mapping. \\
\hline
\textbf{Q4 -- Feedback} & Progress bar, selectable options, result with card and details. \\
\hline
\end{tabular}

\vspace{0.3cm}
\noindent The guide is simple and accessible even for users with 1/5 literacy.
\end{tcolorbox}

\vspace{0.3cm}


\vspace{0.3cm}

% =====================================================================
% TASK 5: CARD VERIFICATION AND TOP-UP
% =====================================================================
\section{Task 5: Card Verification and Top-up}

\begin{tcolorbox}[colback=primaryColor!10,colframe=primaryColor,arc=2mm]
\small
\textbf{User profile:} Liliana (70, low digital literacy) wants to check the status of her TPER card and top it up. Operations available without login from ``Quick Actions''.

\textbf{Starting point:} Login page \texttt{solweb/login.html}. \textbf{Goal:} Verify the card and top it up.
\end{tcolorbox}

\vspace{0.3cm}

\subsection{Step 1: Finding ``Verify Card'' in Quick Actions}

\begin{tcolorbox}[colback=backgroundColor,colframe=borderColor,arc=2mm]
\small
\begin{tabular}{|L{3cm}|L{11cm}|}
\hline
\textbf{User action} & \multicolumn{1}{L{11cm}|}{On the login page, in the ``Quick Actions'' section, clicks ``Verify Card''.} \\
\hline
\textbf{Q1 -- Goal} & The user wants to verify their card. ``Verify Card'' is explicit. \\
\hline
\textbf{Q2 -- Visibility} & The Quick Actions are in a 4-button grid below the login form. ``Verify Card'' is the last one, but clearly visible. \\
\hline
\textbf{Q3 -- Mapping} & ``Verify Card'' $\rightarrow$ ``check my card's status'' is direct. \\
\hline
\textbf{Q4 -- Feedback} & Click $\rightarrow$ navigation to \texttt{verifica.html}. \\
\hline
\end{tabular}

\vspace{0.3cm}
\noindent Clear operation, accessible without login.
\end{tcolorbox}

\vspace{0.3cm}

\subsection{Step 2: Entering the card number and viewing the result}

\begin{tcolorbox}[colback=backgroundColor,colframe=borderColor,arc=2mm]
\small
\begin{tabular}{|L{3cm}|L{11cm}|}
\hline
\textbf{User action} & \multicolumn{1}{L{11cm}|}{On \texttt{verifica.html}, enters the card number and clicks ``Verify''. Sees the simulated results: status, expiry, subscription type, last top-up.} \\
\hline
\textbf{Q1 -- Goal} & The user wants to know the card's status. The form is focused on this. \\
\hline
\textbf{Q2 -- Visibility} & Single ``Card Number'' field + ``Verify'' button. Minimal and clear. \\
\hline
\textbf{Q3 -- Mapping} & ``Enter the number'' $\rightarrow$ ``see the details''. Standard mapping. \\
\hline
\textbf{Q4 -- Feedback} & The results are shown in a card with status icons (green = active, red = expired) and details. Excellent and immediate feedback. \\
\hline
\end{tabular}

\vspace{0.3cm}
\noindent Simple page, clear visual result.
\end{tcolorbox}

\vspace{0.3cm}

\subsection{Step 3: Topping up the card}

\begin{tcolorbox}[colback=backgroundColor,colframe=borderColor,arc=2mm]
\small
\begin{tabular}{|L{3cm}|L{11cm}|}
\hline
\textbf{User action} & \multicolumn{1}{L{11cm}|}{From the login page, clicks ``Top Up Card'' $\rightarrow$ \texttt{ricarica.html}. Enters the card number, chooses an amount from the presets (10\euro, 20\euro, 50\euro, or custom amount), clicks ``Top Up'' $\rightarrow$ payment.} \\
\hline
\textbf{Q1 -- Goal} & The user wants to top up. ``Top Up Card'' is explicit. \\
\hline
\textbf{Q2 -- Visibility} & Clearly visible button in Quick Actions. On the page, clear form with large, clickable preset amounts. \\
\hline
\textbf{Q3 -- Mapping} & ``Choose an amount'' $\rightarrow$ ``Top Up''. Mapping similar to phone top-ups, familiar to many users. \\
\hline
\textbf{Q4 -- Feedback} & After payment, redirect to \texttt{conferma.html} with summary. \\
\hline
\end{tabular}

\vspace{0.3cm}
\noindent Quick flow, preset amounts reduce cognitive effort.

\vspace{0.2cm}
\noindent\textbf{Issue identified:} \textbf{CW-07 (Minor).} It is not clear to the user what happens after the top-up (e.g., ``when will the funds be available on the card?''). A confirmation message with timing information is recommended.

\vspace{0.2cm}
\noindent\textbf{Issue:} \textbf{CW-11 (Minor, Q1).} After verifying the card, the user does not have a direct link to the top-up. They must remember the path back to the Quick Actions. For Liliana, manually retracing the navigation is an unnecessary obstacle.
\end{tcolorbox}

\vspace{0.3cm}


\vspace{0.3cm}

% =====================================================================
% OVERALL SUMMARY
% =====================================================================
\section{Overall Summary}

\begin{tcolorbox}[colback=backgroundColor,colframe=primaryColor,arc=2mm]
\small

\renewcommand{\arraystretch}{1.3}
\begin{tabular}{|L{3.5cm}|C{1.5cm}|C{3cm}|}
\hline
\textbf{Task} & \textbf{Steps} & \textbf{Issues found} \\
\hline
\#1 Guest Renewal & 5 & CW-01, CW-02, CW-08 \\
\hline
\#2 Login Purchase & 4 & CW-03, CW-04 \\
\hline
\#3 Concessions & 3 & CW-05, CW-10 \\
\hline
\#4 Where to Buy & 3 & CW-06, CW-09 \\
\hline
\#5 Verify/Top-up & 3 & CW-07, CW-11 \\
\hline
\rowcolor{primaryColor!10}
\textbf{Total} & \textbf{18} & \textbf{11 issues} \\
\hline
\end{tabular}

\vspace{0.3cm}
\noindent\textbf{Result:} The Cognitive Walkthrough identified \textbf{11 critical issues} out of 18 analyzed steps (4 moderate severity, 7 minor). No step showed critical issues that would prevent task completion. The problems mainly concern communication (ambiguous labels, post-action messages) rather than the overall system architecture.

\end{tcolorbox}

\vspace{0.3cm}

\subsection{Identified Issues}

\begin{tcolorbox}[colback=accentOrange!5,colframe=accentOrange,arc=1mm]
\small
\renewcommand{\arraystretch}{1.2}
\begin{tabular}{|L{1cm}|L{2.5cm}|L{2cm}|L{9cm}|}
\hline
\textbf{ID} & \textbf{Task} & \textbf{Severity} & \multicolumn{1}{L{9cm}|}{\textbf{Description}} \\
\hline
CW-01 & \#1 Renewal & Minor & ``Where to Buy'' is not immediately associated with ``renewal''. Mitigated by the alternative path via ``Subscriptions''. \\
\hline
CW-02 & \#1 Renewal & Minor & Card format validation missing (TPER-XXXXXX format not validated client-side). \\
\hline
CW-03 & \#2 Purchase & \textbf{Medium} & ``And now? What to do after purchase'' section present but poorly visible (below the fold). The user may not see it and board without a valid ticket. \\
\hline
CW-04 & \#2 Purchase & Minor & ``Fiscal Code'' not international. The placeholder ``or Email'' only partially mitigates. \\
\hline
CW-05 & \#3 Concessions & Minor & Missing ``Print summary'' in the wizard results. \\
\hline
CW-06 & \#4 Guide & \textbf{Medium} & ``You cannot buy on tper.it'' message may disorient vulnerable users. A ``Buy now'' one-click option is needed. \\
\hline
CW-07 & \#5 Top-up & Minor & Missing post-top-up timing communication. \\
\hline
CW-08 & \#1 Renewal & \textbf{Medium} & Quick Actions poorly visible below the dominant login form. Risk of abandonment for users with low digital literacy (Q2 Visibility). \\
\hline
CW-09 & \#4 Guide & \textbf{Medium} & Language selector ``IT/EN/AR'' illegible for those who do not know the Latin alphabet. Fatima may not discover the Arabic version (Q2 Visibility). \\
\hline
CW-10 & \#3 Concessions & \textbf{Medium} & ``No concessions'' result perceived as failure. Missing digital alternative to keep the user in the flow (Q4 Feedback). \\
\hline
CW-11 & \#5 Verify & Minor & No direct link from card verification to top-up. Must remember and retrace the navigation (Q1 Goal). \\
\hline
\end{tabular}
\end{tcolorbox}

\vspace{0.3cm}

\subsection{Recommendations}

\begin{tcolorbox}[colback=backgroundColor,colframe=accentGreen,arc=2mm]
\small
Based on the Cognitive Walkthrough results, the following interventions are recommended:

\begin{enumerate}[leftmargin=*]
 \item \textbf{Improve visibility of the post-purchase section (CW-03):} Move the ``And now? What to do after purchase'' section above the fold, in a prominent position relative to the order summary, or replicate it as a confirmation message before the receipt download.

 \item \textbf{Add a ``Buy now'' path (CW-06):} For users with very low digital literacy, add a ``I want to buy a ticket now'' button on the ``Where to Buy'' page that leads directly to an operational channel (e.g., Roger App), bypassing the guide.

 \item \textbf{Add a print summary (CW-05):} Add a ``Print'' button in the concessions wizard results and in document checklists, to support the analog-digital bridge.

 \item \textbf{Card format validation (CW-02):} Add client-side validation for the ``TPER-XXXXXX'' format with immediate feedback.

 \item \textbf{Top-up timing communication (CW-07):} Add a post-top-up confirmation message indicating fund availability timelines.

 \item \textbf{Improve Quick Actions visibility (CW-08):} Move guest operations above the login form or to a more visible position, so that elderly users notice them before confronting the login form.

 \item \textbf{Language selector with native labels (CW-09):} Add the target-language labels alongside the ``IT/EN/AR'' abbreviations (``Italiano'', ``English'', ``Arabo''), so that every user recognizes their own language.

 \item \textbf{Digital alternative to ``no concessions'' (CW-10):} Replace the failure message with ``View all available subscriptions'' or ``Discover generic reductions'', keeping the user in the digital flow.

 \item \textbf{Direct verification $\rightarrow$ top-up link (CW-11):} Add a ``Top up this card'' button on \texttt{verifica.html} that leads directly to \texttt{ricarica.html} with the card number pre-filled.
\end{enumerate}

\vspace{0.3cm}
\noindent\textbf{Priority:} High (CW-03, CW-06, CW-08, CW-09), Medium (CW-05, CW-07, CW-10), Low (CW-01, CW-02, CW-04, CW-11).
\end{tcolorbox}

\vspace{0.3cm}

\subsection*{Name of Redesigned Version}
\begin{tcolorbox}[colback=backgroundColor,colframe=borderColor,arc=2mm]
\small
TPER Redesign --- v2.0 (Interactive Prototype --- Phase D)
\end{tcolorbox}


